\documentclass[10pt,twocolumn]{article}

\usepackage{amsmath,amssymb,amsthm}
\usepackage{bm}
\usepackage{booktabs}
\usepackage{multirow}
\usepackage{graphicx}
\usepackage{xcolor}
\usepackage{algorithm}
\usepackage{algpseudocode}
\usepackage{microtype}
\usepackage{cite}
\usepackage{enumitem}
\usepackage[top=0.75in,bottom=0.72in,left=0.68in,right=0.68in,columnsep=0.22in]{geometry}
\usepackage[colorlinks=true,linkcolor=blue!55!black,citecolor=blue!55!black,urlcolor=blue!55!black]{hyperref}

\renewcommand{\dbltopfraction}{0.90}
\renewcommand{\dblfloatpagefraction}{0.80}
\renewcommand{\topfraction}{0.90}
\renewcommand{\textfraction}{0.08}
\setlength{\textfloatsep}{9pt plus 2pt minus 2pt}
\setlength{\dbltextfloatsep}{8pt plus 2pt minus 2pt}

\definecolor{dtragreen}{RGB}{16,138,67}
\definecolor{abkblue}{RGB}{41,105,216}

\title{\textbf{DTRA: Real-Time Continual 3D Manipulation via Two-Stage Docking and Elastic Snap-Back of Hypertube Rings}}

\author{
  Seungjoon Son \\
  Independent AI \& Robotics Researcher, Daejeon, Republic of Korea \\
  \texttt{\href{https://github.com/5Color/dtra}{https://github.com/5Color/dtra}}
}

\date{November 2026}

\begin{document}

\maketitle

\begin{abstract}
Task-ID-free continual visuomotor manipulation requires simultaneously retaining prior skills without catastrophic forgetting, disambiguating visually similar reset configurations, and maintaining constant inference latency as the learned task repertoire $K$ grows. Existing slot-isolated architectures such as ABK-VLA v1.1.3~\cite{son2026abkvla} guarantee zero forgetting on small task suites, yet reveal three structural bottlenecks when scaled to 8 sequential 3D manipulation tasks: (i) $O(K)$ sequential residual chaining causes CPU chunk latency to grow linearly ($5.08\,\text{ms} \to 27.48\,\text{ms}$), (ii) fixed $t=0$ Difference-of-Gaussians (DoG) visual keys collapse on visually symmetric initial scenes (minimum cosine margin $0.0082$, inducing $8.57\%$ forgetting under Blend routing), and (iii) high-frequency edge filters degrade sharply under camera noise and brightness shifts ($4.7\%$ success rate). We introduce \textbf{DTRA} (Dynamic Topological Ring Architecture), a first-principles non-destructive policy framework that models procedural skills as closed first-order elliptical hypertube rings embedded on the unit hypersphere $\mathbb{S}^{63}$ ($D_{\text{topo}}=64$). Each tube possesses a learnable radius $r_i$, bounded elastic compliance $\alpha_i \in (0,1)$, and multimodal visual-language magnetic key $\mathbf{k}_i \in \mathbb{R}^{128}$. Driven by a unified magnetic attraction force integrating intra-episode inertia and command-compatible myelinated associative coupling $M_{ij}$, incoming queries dock onto sparsely gated active tubes ($O(|A_t|) \approx O(1)$) via closed-form $\operatorname{atan2}$ phase projection across a two-stage Coarse-to-Fine hierarchy, followed by an instantaneous post-inference elastic snap-back ($\Delta\mathbf{z}_i \leftarrow \mathbf{0}$) that enforces exact topological preservation. Across six Step~1 dynamical verification checks, DTRA achieves exact $0.0$ snap-back error, $25.0\%$ sparse eval activation, $0.0$ shape drift under structural outgrowth, and $+0.1648$ command-gated associative recall with $0.0$ misfire on incompatible queries. On an 8-task 3D Meta-World continual benchmark, our DTRA~(v2.1) implementation attains \textbf{90.0\% Avg SR} with \textbf{0.0\% FR} and \textbf{0.0\% BWT} (outperforming ABK-VLA Blend at $80.0\%$ SR / $8.57\%$ FR and Argmax at $82.5\%$ SR), maintains constant $O(1)$ CPU inference speed ($10.98\,\text{ms}$ at $K=2$ vs.\ $10.28\,\text{ms}$ at $K=8$), widens the average inter-task resonance margin by \textbf{13.9$\times$} ($0.2752$ vs.\ $0.0198$), and achieves \textbf{9$\times$ higher severe OOD robustness} ($42.2\%$ vs.\ $4.7\%$).
\end{abstract}

% ==========================================================================
\section{Introduction}
% ==========================================================================

Deploying generalist visuomotor and Vision-Language-Action (VLA) policies~\cite{kim2024openvla,zhao2023act,chi2023diffusion} in sequential open-ended environments requires learning new manipulation skills over time without degrading previously mastered behaviors~\cite{wolczyk2021continualworld,liu2023libero}. Classical weight-regularization methods such as Elastic Weight Consolidation (EWC)~\cite{kirkpatrick2017ewc} and Riemannian Walk~\cite{chaudhry2018riemannian} protect parameters quadratically but struggle under severe plasticity--stability tension across heterogeneous continuous control tasks. Parameter-isolation architectures such as Progressive Neural Networks~\cite{rusu2016progressive}, PackNet~\cite{mallya2018packnet}, and prompt-routed pools~\cite{wang2022l2p} preserve prior skills through frozen sub-networks, while ABK-VLA v1.1.3~\cite{son2026abkvla} unified frozen cubic B-Spline trajectory compression~\cite{paraschos2013promp} with task-ID-free $t=0$ visual prototype routing on 4-task 2D sketch and 3D Meta-World suites.

However, scaling ABK-VLA v1.1.3 from $K=4$ to $K=8$ sequential 3D Meta-World manipulation tasks~\cite{yu2020metaworld} exposes three fundamental limitations:
\begin{enumerate}[leftmargin=*,itemsep=2pt]
  \item \textbf{$O(K)$ Sequential Residual Latency Scaling:} Because ABK-VLA accumulates delta control points through sequential weighted slot summations across $s=0,\dots,K-1$, per-chunk inference runtime on edge CPUs grows from $5.08\,\text{ms}$ ($K=1$) to $27.48\,\text{ms}$ ($K=8$).
  \item \textbf{$t=0$ Visual Key Collision on Symmetric Scenes:} Fixed $15\times 15$ Difference-of-Gaussians (DoG) visual prototypes yield razor-thin inter-task cosine margins ($0.0082$ between \texttt{drawer-open-v3} and \texttt{drawer-close-v3}), triggering $8.57\%$ forgetting under soft Blend routing.
  \item \textbf{Sensor Noise Fragility:} High-pass DoG edge filtering amplifies high-frequency camera corruption, collapsing autonomous routing accuracy to $40.6\%$ and manipulation success rate to $4.7\%$ under severe Gaussian noise ($\sigma=0.20$) and brightness shifts.
\end{enumerate}

To overcome these bottlenecks from first principles, we propose \textbf{DTRA} (\emph{Dynamic Topological Ring Architecture}), modeling motor knowledge not as static flat vectors or chained residual heads, but as \textbf{3D-intuitive closed elliptical hypertube rings} embedded on a shared 64-dimensional unit hypersphere $\mathbb{S}^{63} \subset \mathbb{R}^{64}$ (Figure~\ref{fig:step1_dynamics}). Our core contributions are summarized as follows:
\begin{itemize}[leftmargin=*,itemsep=2pt]
  \item \textbf{First-Principles Unit-Hypersphere Hypertube Formulation:} We formalize each skill primitive as a closed first-order harmonic ring $R_i^{(0)}(\theta)$ on $\mathbb{S}^{63}$ parameterized by learnable thickness $r_i$, bounded elastic compliance $\alpha_i = \sigma(\omega_i) \in (0,1)$, and adaptive magnetic gate $\tau_i = \operatorname{softplus}(\beta_i)$. Replacing discretized phase grid search with an exact closed-form $\operatorname{atan2}$ projection yields analytical nearest-point docking under a unified spherical chord distance $d_{\text{chord}}$.
  \item \textbf{Unified Magnetic Attraction, Command-Compatible Myelination \& Instantaneous Snap-Back:} Incoming multimodal queries attract only active tubes above $\tau_i$ via intra-episode cross-chunk inertia ($\eta\, o_{i,t-1}$) and command-gated Hebbian myelination ($\sum_{j \neq i} M_{ij} a_j(t-1)\max(0,\cos(\mathbf{x}_t, \mathbf{k}_i))$), transiently deforming tubes during inference before snapping back ($\Delta\mathbf{z}_i \leftarrow \mathbf{0}$) to guarantee $0.0$ long-term drift and $0\%$ forgetting.
  \item \textbf{Multimodal Coarse-to-Fine 3D Manipulation \& $O(1)$ CPU Execution:} Fusing mid-level Conv4 spatial features with an offline lightweight language encoder cleanly separates visually identical table setups ($13.9\times$ wider mean cosine margin, $0.2752$ vs.\ $0.0198$), while two-stage Coarse-to-Fine active-tube slicing delivers $O(1)$ CPU latency ($10.28\,\text{ms}$ at $K=8$), \textbf{90.0\% Avg SR}, \textbf{0.0\% FR}, and \textbf{9$\times$ severe OOD SR} ($42.2\%$ vs.\ $4.7\%$).
\end{itemize}

% ==========================================================================
\section{Related Work}
% ==========================================================================

\paragraph{Continual Learning and Modular Isolation.}
Continual learning algorithms balance knowledge retention against forward adaptability~\cite{chaudhry2018riemannian}. Regularization methods such as EWC~\cite{kirkpatrick2017ewc} and RWalk~\cite{chaudhry2018riemannian} constrain weights near prior task optima, whereas episodic memory approaches like GEM~\cite{lopezpaz2017gem} project gradients to prevent retroactive loss increases. Parameter-isolation architectures eliminate interference explicitly via dedicated task columns~\cite{rusu2016progressive}, iterative binary weight masks~\cite{mallya2018packnet}, or key--query prompt pools~\cite{wang2022l2p}. DTRA inherits structural isolation and task-ID-free routing while replacing flat prompt tokens with geometrically bounded, elastically deformable hypertube rings equipped with explicit pairwise repulsion margins on $\mathbb{S}^{63}$.

\paragraph{Continual Robot Manipulation and Action Chunking.}
Standard benchmarks such as Meta-World~\cite{yu2020metaworld}, Continual World~\cite{wolczyk2021continualworld}, and LIBERO~\cite{liu2023libero} highlight that sequential fine-tuning degrades procedural manipulation skills unless architectures structurally separate task representations. Meanwhile, Action Chunking with Transformers (ACT)~\cite{zhao2023act}, Diffusion Policy~\cite{chi2023diffusion}, and generalist VLA models such as OpenVLA~\cite{kim2024openvla} demonstrate that predicting temporally extended motor chunks stabilizes closed-loop visuomotor control~\cite{levine2016visuomotor}. Compressing action chunks through movement primitives~\cite{paraschos2013promp} or clamped cubic B-Spline control points~\cite{son2026abkvla} shrinks head dimensionality. DTRA preserves ABK-VLA's frozen B-Spline action interface while eliminating its $O(K)$ residual sum and DoG key collisions.

% ==========================================================================
\section{Methodology: DTRA}
% ==========================================================================

\subsection{Multimodal Magnetic Key Construction}
To disambiguate tasks that share identical initial table geometry (e.g., pressing a desk button vs.\ pressing a coffee-machine button, or opening vs.\ closing a drawer), DTRA replaces manual high-pass $15\times 15$ DoG filtering with a $D_{\text{key}}=128$ multimodal magnetic query vector $\mathbf{x}_t \in \mathbb{S}^{127}$ combining illumination-resilient spatial features and natural-language instructions.

Given an RGB observation $\mathbf{I}_t \in \mathbb{R}^{3 \times 128 \times 128}$, we extract the mid-level Conv4 feature map from the spatial encoder (frozen after Task~0 in continual 3D learning), apply adaptive average pooling to $(4\times 4)$ spatial grid resolution (flattening to $\mathbb{R}^{4096}$), and project it via a fixed semi-orthogonal QR matrix $\mathbf{Q}_{\text{vis}} \in \mathbb{R}^{4096 \times 64}$ followed by $L_2$ normalization to obtain $\mathbf{v}_t \in \mathbb{S}^{63}$. Simultaneously, the task command $L_{\text{task}}$ is encoded by an offline lightweight deterministic unigram/character-3-gram hash projection (\texttt{LightweightTextEncoder}, $512 \to 64$ normalized via QR basis) to yield $\boldsymbol{\ell}_{\text{task}} \in \mathbb{S}^{63}$. Both modalities are fused cleanly onto the unit hypersphere:
\begin{equation}
  \mathbf{x}_t = \operatorname{norm}\!\Bigl(\bigl[\,0.55\,\mathbf{v}_t \;;\; 0.835\,\boldsymbol{\ell}_{\text{task}}\,\bigr]\Bigr) \in \mathbb{S}^{127},
\end{equation}
where $\operatorname{norm}(\mathbf{u}) = \mathbf{u} / \max(\|\mathbf{u}\|_2, \epsilon)$. Weighting language slightly above low-level visual variation ($0.835$ vs.\ $0.55$) ensures deterministic separation under camera noise while retaining frame-level visual grounding inside $\mathbf{v}_t$.

\subsection{First-Order Hypertube Geometry on \texorpdfstring{$\mathbb{S}^{63}$}{S\textasciicircum 63}}
Every skill primitive $i$ lives on a shared $D_{\text{topo}}=64$ unit hypersphere $\mathbb{S}^{63} \subset \mathbb{R}^{64}$ via a shared linear projection $P_{\text{topo}} \in \mathbb{R}^{64 \times 128}$. The idle spine (\emph{base ring}) $\mathbf{z}_i^{(0)}(\theta)$ of tube $R_i$ is parameterized by the coordinate-synchronized key center and a closed \textbf{first-order elliptical Fourier basis} $B_{\text{ring}}(\theta) = [1,\cos\theta,\sin\theta]^\top$ with trainable harmonic matrix $C_i = [\mathbf{c}_{i,\text{dc}}, \mathbf{c}_{i,\cos}, \mathbf{c}_{i,\sin}]^\top \in \mathbb{R}^{3 \times 64}$ (initialized with orthogonal planar vectors $\mathbf{c}_{i,\cos} \perp \mathbf{c}_{i,\sin}$):
\begin{equation}
\begin{aligned}
  \mathbf{z}_i^{(0)}(\theta) = \operatorname{norm}\!\bigl(
    & P_{\text{topo}}\hat{\mathbf{k}}_i + \mathbf{s}_i + \mathbf{c}_{i,\text{dc}} \\
    & + \cos\theta\,\mathbf{c}_{i,\cos} + \sin\theta\,\mathbf{c}_{i,\sin}
  \bigr) \in \mathbb{S}^{63},
\end{aligned}
\label{eq:base_ring}
\end{equation}
where $\hat{\mathbf{k}}_i = \operatorname{norm}(\mathbf{k}_i)$ and $\mathbf{s}_i \in \mathbb{R}^{64}$ denotes optional structural outgrowth displacement ($\mathbf{s}_i=\mathbf{0}$ prior to outgrowth). Each tube carries three physical scalar parameters:
\begin{equation}
\begin{aligned}
  r_i &= \operatorname{softplus}(\rho_i) + 10^{-3}, \\
  \alpha_i &= \sigma(\omega_i) \in (0,1), \qquad
  \tau_i = \operatorname{softplus}(\beta_i),
\end{aligned}
\end{equation}
governing respectively the tube's effective attraction thickness $r_i > 0$, bounded elastic compliance (\emph{softness}) $\alpha_i \in (0,1)$ avoiding division-by-zero stiffness singularities, and adaptive magnetic gate bias $\tau_i \ge 0$. All pairwise metrics on $\mathbb{S}^{63}$ use the unified \textbf{spherical chord distance}:
\begin{equation}
\begin{aligned}
  d_{\text{chord}}(\mathbf{u}, \mathbf{v})
  &= \|\mathbf{u} - \mathbf{v}\|_2 \\
  &= \sqrt{2\bigl(1 - \operatorname{clamp}(\langle \mathbf{u}, \mathbf{v}\rangle, -1, 1)\bigr) + \epsilon}.
\end{aligned}
\end{equation}

\paragraph{Closed-Form Analytical $\operatorname{atan2}$ Projection.}
Given an incoming topological query $\mathbf{c}_t = \operatorname{norm}(P_{\text{topo}}\mathbf{x}_t) \in \mathbb{S}^{63}$, projecting $\mathbf{c}_t$ onto the planar ellipse spanned by $(\mathbf{c}_{i,\cos}, \mathbf{c}_{i,\sin})$ maximizes the alignment $\langle \mathbf{c}_t, \cos\theta\,\mathbf{c}_{i,\cos} + \sin\theta\,\mathbf{c}_{i,\sin}\rangle$ via a single exact closed-form phase angle:
\begin{equation}
  \theta_i^*(\mathbf{c}_t) = \operatorname{atan2}\!\Bigl(\langle \mathbf{c}_t, \mathbf{c}_{i,\sin}\rangle,\; \langle \mathbf{c}_t, \mathbf{c}_{i,\cos}\rangle\Bigr),
\end{equation}
yielding the exact nearest anchor $\mathbf{z}_{i,0}^* = \mathbf{z}_i^{(0)}(\theta_i^*)$ in $O(1)$ independent of phase sample count $M_\theta$.

\subsection{Unified Magnetic Force, Myelination \& Snap-Back}
Across sequential chunks $t=1,2,\dots$ within an episode, DTRA evaluates the unified magnetic attraction force across normalized keys:
\begin{equation}
\begin{aligned}
  F_{\text{mag}}^{(i)}(t)
  &= \cos(\mathbf{x}_t, \mathbf{k}_i) + \eta\, o_{i, t-1} \\
  &\quad + \sum_{j \neq i} M_{ij}\, a_j(t-1)\,\max\!\bigl(0, \cos(\mathbf{x}_t, \mathbf{k}_i)\bigr),
\end{aligned}
\label{eq:fmag}
\end{equation}
where $\eta=0.25$ propagates cross-chunk topological continuity from the previous overlap $o_{i,t-1}$ and $M_{ij} \in [0, \eta_{\text{myelin}}]$ ($\eta_{\text{myelin}}=0.35$) denotes symmetric myelinated associative conductance learned from co-docking history. Crucially, the multiplicative compatibility factor $\max(0, \cos(\mathbf{x}_t, \mathbf{k}_i))$ suppresses associative hallucination whenever the active command $\mathbf{x}_t$ opposes or is orthogonal to skill $\mathbf{k}_i$.

Defining $\Delta F_i(t) = F_{\text{mag}}^{(i)}(t) - \tau_i$, an adaptive magnetic gate prunes inactive rings:
\begin{equation}
  a_i(t) =
  \begin{cases}
    \operatorname{ReLU}(\Delta F_i) + \varepsilon_{\text{tail}}\,\sigma(\gamma_{\text{tail}}\Delta F_i), & \text{Train}, \\
    \operatorname{ReLU}(\Delta F_i), & \text{Eval},
  \end{cases}
  \label{eq:gate}
\end{equation}
with soft exploration tail ($\varepsilon_{\text{tail}}=0.05, \gamma_{\text{tail}}=4.0$) preserving gradients during training while slicing strictly to the sparse active set $A_t = \{i : a_i(t) > 0\}$ during evaluation (falling back to the top-$1$ clamped margin $\max(F_{\text{mag}}-\tau+0.25, 0.05)$ in single-tube 3D inference if all raw margins are non-positive).

\paragraph{Elastic Deformation and Post-Inference Snap-Back.}
Only active tubes $i \in A_t$ bend smoothly toward $\mathbf{c}_t$ via a localized phase kernel $\psi_i(\theta) = \bigl(\tfrac{1 + \cos(\theta - \theta_i^*)}{2}\bigr)^2$:
\begin{equation}
  \Delta \mathbf{z}_i^* = \alpha_i\,a_i(t)\bigl(\mathbf{c}_t - \mathbf{z}_{i,0}^*\bigr), \quad
  \tilde{\mathbf{z}}_i^* = \operatorname{norm}\!\bigl(\mathbf{z}_{i,0}^* + \Delta\mathbf{z}_i^*\bigr) \in \mathbb{S}^{63},
\end{equation}
with full-ring state buffer $\Delta\mathbf{z}_i(\theta) = \Delta\mathbf{z}_i^*\,\psi_i(\theta)$. The \textbf{Soft Tube Overlap} then measures whether $\mathbf{c}_t$ lies inside the deformed tube radius $r_i$:
\begin{equation}
  o_{i,t} = \tanh\!\bigl(g_{\text{gate}}\,a_i(t)\bigr)\,\sigma\!\left(\frac{r_i - d_{\text{chord}}(\mathbf{c}_t, \tilde{\mathbf{z}}_i^*)}{T_{\text{ov}}}\right),
\end{equation}
($g_{\text{gate}}=4.0$; $T_{\text{ov}}=0.14$ in Step~1, $0.16$ in 3D). Immediately after action synthesis, \texttt{snap\_back()} clears $\Delta\mathbf{z}_i(\theta) \leftarrow \mathbf{0}$ so every stored idle orbit returns identically to $\mathbf{z}_i^{(0)}$ ($\|\tilde{\mathbf{z}}_i - \mathbf{z}_i^{(0)}\|_\infty = 0.0$); at episode boundaries, \texttt{reset()} zeros $o_{i,0}=0$, $a_i(0)=0$, and $\Delta\mathbf{z}_i=\mathbf{0}$.

\subsection{Two-Stage Docking \& Cubic B-Spline Action Synthesis}
DTRA stacks two hyperspherical ring stages sharing $P_{\text{topo}}$:
\begin{enumerate}[leftmargin=*,itemsep=2pt]
  \item \textbf{Stage~1 (Coarse Semantic Primitive):} Thick, compliant rings ($r_{\text{init}}=0.40$--$0.45$, $\alpha_{\text{init}}=0.78$) compute normalized overlap weights $w_i^c = o_{i,t}^c / (\sum_{j \in A_t^c} o_{j,t}^c + 10^{-4})$ and aggregated representation $(\mathbf{z}_{\text{dock}}^c, \mathbf{h}_{\text{avg}}^c)$ (plus second-order interference pooling $\mathbf{h}_{\text{int}}$ in the full prototype). A \texttt{CoarseAnchorHead} emits base cubic B-Spline control points $\mathbf{P}_{\text{anchor}} \in (0,1)^{N_k \times d_a}$ and a task context shift $\mathbf{c}_{\text{ctx}} \in (-1,1)^{128}$, redirecting Stage~2 queries via $\mathbf{x}_t^{\text{fine}} = \operatorname{norm}(\mathbf{x}_t + 0.35\,\mathbf{c}_{\text{ctx}})$ in 3D ($\operatorname{LN}(\mathbf{x}_t + 0.5\,\mathbf{c}_{\text{ctx}})$ in Step~1).
  \item \textbf{Stage~2 (Fine Motor Precision):} Thin, stiff rings ($r_{\text{init}}=0.18$--$0.25$, $\alpha_{\text{init}}=0.24$) refine precision docking to output $(\mathbf{z}_{\text{dock}}^f, \mathbf{h}_{\text{avg}}^f, o_{i,t}^f)$. An $M=6$ Winner-Takes-All (WTA) residual head predicts multimodal control-point offsets $\Delta\mathbf{P}_m$ and selector logits, scaling the winning residual by overlap resonance ($0.55\,o_{i,t}^f\tanh(\Delta\mathbf{P}_{m^*})$ in 3D). Final control points $\mathbf{P}^* = \operatorname{clamp}(\mathbf{P}_{\text{anchor}} + \Delta\mathbf{P}^*, 0, 1)$ are decoded by a frozen clamped cubic B-Spline matrix~\cite{son2026abkvla} from $N_k=6$ knots to a $H=16$-step continuous trajectory chunk (executed open-loop for $8$ steps per replan).
\end{enumerate}

\subsection{Two-Term Regularization \& Non-Destructive Growth}
To prevent unbounded tube inflation (\emph{lazy expansion}) and inter-tube manifold collisions without redundant hyper-parameters, DTRA employs a first-principles \textbf{two-term topological loss}:
\begin{equation}
  \mathcal{L}_{\text{topo}}
  = \lambda_r \sum_{i} r_i
  + \lambda_{\text{rep}} \sum_{i < j} \operatorname{ReLU}\!\bigl(r_i + r_j + \delta - d_{ij}^{(0)}\bigr)^2,
  \label{eq:l_topo}
\end{equation}
where $d_{ij}^{(0)} \equiv d_{\text{chord}}(R_i^{(0)}, R_j^{(0)}) = \min_{\theta,\phi} d_{\text{chord}}(\mathbf{z}_i^{(0)}(\theta), \mathbf{z}_j^{(0)}(\phi))$ measures minimum base-ring separation over $M_\theta$ sampled phases ($\delta=0.06$--$0.08$).

\paragraph{Hebbian Edge Myelination \& Structural Outgrowth.}
Co-activated tube pairs $(i,j)$ accumulate symmetric coactivation traces $C_{ij}$, forming conductive shortcuts $M_{ij} = \eta_{\text{myelin}}\tanh(C_{ij})$ while proportionally shifting only center offsets $\mathbf{s}_i, \mathbf{s}_j$ outside the repulsion safety floor $r_i+r_j+\delta$ and holding $(C_i, r_i, \alpha_i)$ strictly invariant.

\paragraph{Non-Destructive Task Promotion in Continual 3D Control.}
Upon encountering a new task in continual learning, the shared trunk (\texttt{encoder}, \texttt{state\_net}, $P_{\text{topo}}$) and all prior Coarse/Fine nodes remain frozen after Task~0. If mean query cosine alignment with existing nodes falls below $\theta_c=0.74$ (Coarse) or $\theta_f=0.96$ (Fine), a new \texttt{HypertubeNode} is appended and trained in isolation under a $6$-mode WTA B-Spline objective plus weighted $\mathcal{L}_{\text{topo}}$ against frozen base rings. Algorithm~\ref{alg:dtra_infer} summarizes closed-loop chunk execution.

\begin{algorithm}[t]
\caption{DTRA \texorpdfstring{$O(1)$}{O(1)} Closed-Loop Inference \& Snap-Back}
\label{alg:dtra_infer}
\begin{algorithmic}[1]
\Require Episode start: call \texttt{reset()} ($o_{i,0} \leftarrow 0, a_i(0) \leftarrow 0, \Delta\mathbf{z}_i \leftarrow \mathbf{0}$)
\For{replanning step $t = 0, 8, 16, \dots$}
  \State Extract $\mathbf{x}_t = \operatorname{norm}([0.55\,\mathbf{v}_t(\mathbf{I}_t);\; 0.835\,\boldsymbol{\ell}_{\text{task}}])$
  \State Compute $F_{\text{mag}}^c(t)$ across Coarse bank (Eq.~\ref{eq:fmag}); select sparse active $A_t^c$ ($O(|A_t^c|)$)
  \State Closed-form $\theta_i^* \leftarrow \operatorname{atan2}$; deform $\tilde{\mathbf{z}}_i^*$ and compute $o_{i,t}^c$ for $i \in A_t^c$
  \State $(\mathbf{P}_{\text{anchor}}, \mathbf{c}_{\text{ctx}}) \leftarrow \texttt{CoarseHead}([\mathbf{c}_t, \mathbf{z}_{\text{dock}}^c, \mathbf{h}_{\text{avg}}^c])$
  \State $\mathbf{x}_t^{\text{fine}} \leftarrow \operatorname{norm}(\mathbf{x}_t + 0.35\,\mathbf{c}_{\text{ctx}})$; select sparse active Fine $A_t^f$ ($O(|A_t^f|)$)
  \State Compute Fine $\tilde{\mathbf{z}}_j^*, o_{j,t}^f$; decode $\Delta\mathbf{P}_{m^*}$ via WTA Fine head
  \State $\mathbf{P}^* \leftarrow \operatorname{clamp}(\mathbf{P}_{\text{anchor}} + \Delta\mathbf{P}^*, 0, 1)$; chunk $\hat{\mathbf{Y}}_t \leftarrow \mathbf{B}_{\text{cubic}}\,\mathbf{P}^*$
  \State Execute $\hat{\mathbf{Y}}_{t, 1:8}$ in environment; update buffers $(o_{t}, a_{t})$
  \State \textbf{Snap-back:} call \texttt{snap\_back()} ($\Delta\mathbf{z}_i \leftarrow \mathbf{0}, \forall i$)
\EndFor
\end{algorithmic}
\end{algorithm}

% ==========================================================================
\section{Step 1: Dynamical Verification on \texorpdfstring{$\mathbb{S}^{63}$}{S\textasciicircum 63}}
% ==========================================================================

\begin{figure*}[t]
  \centering
  \includegraphics[width=0.74\textwidth]{../outputs/dtra_step1_dynamics_clean.png}
  \caption{\textbf{DTRA~(v2.1) Step~1 topological dynamics verification on $\mathbb{S}^{63}$ (PCA 3D projection).}
  \textbf{Panel~1:} Idle manifold after 200 optimization steps showing thick, compliant Coarse rings ($r=0.18, \alpha=0.82$) and thin, stiff Fine rings ($r=0.08, \alpha=0.25$).
  \textbf{Panel~2:} Close-up of sparse evaluation docking ($1/4$ Coarse, $2/8$ Fine active; $O(|A_t|)\approx O(1)$), where closed-form $\operatorname{atan2}$ bends only active rings locally toward the query $\mathbf{c}(\mathbf{x}_t)$.
  \textbf{Panel~3:} Hebbian myelination ($M_{ij} \to 0.350$) and shape-preserving structural outgrowth over 12 co-docking steps: pairwise chord distance converges safely outside the repulsion floor ($1.6385 \to 1.3422$, floor $0.2073$) with $0.0$ shape drift and $+0.165$ ($+0.1648$) compatible vs.\ $0.0$ incompatible command-gated recall.
  \textbf{Panel~4:} Log-scale convergence under two-term $\mathcal{L}_{\text{topo}}$ ($\text{Traj MSE}: 0.05399 \to 0.00002$) alongside exact post-inference snap-back ($L_\infty = 0.0$).}
  \label{fig:step1_dynamics}
\end{figure*}

Before integrating hypertube layers into continuous 3D robot control, we validate all mathematical invariants in a standalone prototype simulation (\texttt{TwoStageDTRA}: $D_{\text{key}}=128, D_{\text{topo}}=64, K_{\text{coarse}}=4, K_{\text{fine}}=8, M_\theta=64, N_k=6, H=16$, trained for 200 AdamW steps over 8 synthetic planar strokes). As summarized in Table~\ref{tab:step1_checks} and visualized in Figure~\ref{fig:step1_dynamics}, all \textbf{six core unit checks} pass without numerical exceptions:

\begin{table}[t]
  \centering
  \caption{Step~1 quantitative unit verification summary on $\mathbb{S}^{63}$ (\texttt{TwoStageDTRA}).}
  \label{tab:step1_checks}
  \small
  \resizebox{\columnwidth}{!}{
  \begin{tabular}{ll}
    \toprule
    \textbf{Verification Check} & \textbf{Measured Quantitative Outcome} \\
    \midrule
    \textbf{1. Gradient Flow} & Coarse $\|\nabla_{\rho,\omega,\beta}\| = (9.90\!\times\!10^{-3}, 1.78\!\times\!10^{-6}, 7.88\!\times\!10^{-5})$; \\
    (Zero NaN/Inf) & Fine $(3.73\!\times\!10^{-2}, 1.04\!\times\!10^{-6}, 8.38\!\times\!10^{-5})$; $\|\nabla P_{\text{topo}}\| = 3.73\!\times\!10^{-3}$ \\
    \textbf{2. Elastic Snap-Back} & Active $L_\infty$ deform: Coarse $0.094370$, Fine $0.006922$; \\
    ($\Delta\mathbf{z}_i \leftarrow \mathbf{0}$) & Post-\texttt{snap\_back()} $\|\tilde{\mathbf{z}}_i - \mathbf{z}_i^{(0)}\|_\infty = \mathbf{0.0}$ (both stages) \\
    \textbf{3. Episode \texttt{reset()}} & Inertia at $t=1$: $(0.2581, 0.0100)$; $F_{\text{mag}}$ boost $t=2$: $(+0.0645, +0.0025)$; \\
    (Cross-chunk state) & Post-\texttt{reset()} buffer $= \mathbf{0.0}$, boost $= \mathbf{0.0}$ \\
    \textbf{4. Sparse Eval Pruning} & Train (soft tail): Coarse $4/4$ ($100.0\%$), Fine $8/8$ ($100.0\%$); \\
    ($O(|A_t|)$ active set) & Eval (pruned): Coarse $\mathbf{1/4}$ ($\mathbf{25.0\%}$), Fine $\mathbf{2/8}$ ($\mathbf{25.0\%}$) \\
    \textbf{5. $\mathcal{L}_{\text{topo}}$ Convergence} & Total Loss $1.2766 \to 0.0099$; Traj MSE $0.05399 \to \mathbf{0.00002}$; \\
    (200 AdamW steps) & $\mathcal{L}_{\text{topo}}: 0.1399 \to 0.0640$; Coarse $(r=0.1832, \alpha=0.8225)$; \\
    & Fine $(r=0.0829, \alpha=0.2488)$; min margins $(+0.3828, +0.4152)$ \\
    \textbf{6. Edge Myelination} & Drift $(\text{shape}, r, \alpha) = (\mathbf{0.0}, 0.0, 0.0)$; $d_{\text{chord}}: 1.6385 \to 1.3422$ \\
    \& Outgrowth & (floor $0.2073, M_{ij}=0.3498$); Compat $\mathbf{+0.1648}$ vs.\ Incompat $\mathbf{0.0}$ \\
    \bottomrule
  \end{tabular}
  }
\end{table}

\begin{itemize}[leftmargin=*,itemsep=2pt]
  \item \textbf{Role Differentiation Emerges Naturally:} Under two-term $\mathcal{L}_{\text{topo}}$ ($\lambda_r^c=0.015, \lambda_r^f=0.080, \lambda_{\text{rep}}=1.60$), Coarse tubes converge to wider, more compliant attractors ($r=0.1832, \alpha=0.8225$) while Fine tubes self-compress into narrow, stiff precision rings ($r=0.0829, \alpha=0.2488$), both maintaining strictly positive separation margins ($+0.3828$ and $+0.4152$).
  \item \textbf{Zero-Misfire Associative Myelination:} Across 12 co-docking updates, $M_{ij}$ reaches $0.3498$ and draws coupled ring centers closer ($1.6385 \to 1.3422$, well outside the $0.2073$ safety floor) with identically $0.0$ distortion of internal Fourier coefficients, radii, or compliance, while boosting compatible queries by $+0.1648$ and blocking incompatible queries to $0.0$.
\end{itemize}

% ==========================================================================
\section{8-Task 3D Meta-World Benchmark}
% ==========================================================================

\subsection{Experimental Protocol}
We expand our prior 4-task 3D Meta-World evaluation~\cite{son2026abkvla} into an \textbf{8-task continual manipulation sequence} ($K=8$, MuJoCo \texttt{ML1} v3, \texttt{corner} RGB camera $128\times 128$, 4-DoF proprioceptive state $\to H=16$ chunks with $N_k=6$ knots, \texttt{exec\_steps=8}):
\begin{center}
  \footnotesize
  $T_0$: \texttt{button-press-v3}, $T_1$: \texttt{door-open-v3}, $T_2$: \texttt{window-open-v3}, $T_3$: \texttt{faucet-open-v3}, \\
  $T_4$: \texttt{drawer-open-v3}, $T_5$: \texttt{drawer-close-v3}, $T_6$: \texttt{handle-press-v3}, $T_7$: \texttt{coffee-button-v3}.
\end{center}
Each task supplies $30$ expert demonstrations and is trained for $55$ epochs (AdamW, $\text{lr}=10^{-3}$, CosineAnnealingLR, \texttt{seed=42}), evaluated sequentially on $15$ held-out initializations per task. We directly compare \textbf{DTRA~(v2.1)} ($5,248,288$ params, allocating $8$ Coarse + $8$ Fine tubes via novelty thresholds $\theta_c=0.74, \theta_f=0.96$) against \textbf{v1.1.3 Autonomous 3D ABK-VLA} ($2,943,031$ params, expanding $8$ residual slots under DoG threshold $0.996$) in both \textbf{Blend} and \textbf{Argmax} modes. Evaluation covers continual manipulation retention (\textbf{Avg SR}, Forgetting Rate \textbf{FR}, Backward Transfer \textbf{BWT}~\cite{lopezpaz2017gem,chaudhry2018riemannian}), $K=1 \dots 8$ latency scaling on isolated CPU and Intel XPU, inter-task cosine separation margins, and an 8-episode-per-task 3D OOD stress test under additive RGB Gaussian noise plus brightness affine shifts: Mild ($\sigma=0.10, 1.15\times + 0.08$) and Severe ($\sigma=0.20, 1.25\times + 0.12$).

\subsection{Quantitative Continual Manipulation Performance}

\begin{table*}[t]
  \centering
  \caption{Final 8-task 3D Meta-World continual manipulation comparison ($R_{i,7}$ per-task success rate \% after training all $K=8$ tasks, $n_{\text{eval}}=15$ held-out episodes per task).}
  \label{tab:mw_results}
  \footnotesize
  \setlength{\tabcolsep}{4.0pt}
  \begin{tabular}{lcccccccccccc}
    \toprule
    \textbf{Method} & \textbf{Params} & $T_0$ & $T_1$ & $T_2$ & $T_3$ & $T_4$ & $T_5$ & $T_6$ & $T_7$ & \textbf{Avg SR $\uparrow$} & \textbf{FR $\downarrow$} & \textbf{BWT $\uparrow$} \\
    \midrule
    v1.1.3 ABK (Blend)  & 2.94M & \textbf{100.0} & \textbf{93.3} & \textbf{73.3} & \textbf{100.0} & 53.3 & 73.3 & 93.3 & 53.3 & 80.0\% & 8.57\% & $-5.71\%$ \\
    v1.1.3 ABK (Argmax) & 2.94M & \textbf{100.0} & 86.7 & 66.7 & \textbf{100.0} & 53.3 & 80.0 & 93.3 & 80.0 & 82.5\% & \textbf{0.00\%} & \textbf{0.00\%} \\
    \textbf{DTRA~(v2.1, Ours)} & 5.25M & 93.3 & 80.0 & 66.7 & \textbf{100.0} & \textbf{93.3} & \textbf{100.0} & \textbf{100.0} & \textbf{86.7} & \textbf{90.0\%} & \textbf{0.00\%} & \textbf{0.00\%} \\
    \bottomrule
  \end{tabular}
\end{table*}

\begin{figure*}[t]
  \centering
  \includegraphics[width=0.75\textwidth]{../outputs/dtra_3d_metaworld_benchmark.png}
  \caption{\textbf{8-Task 3D Meta-World continual benchmark comparing DTRA~(v2.1) against ABK-VLA v1.1.3.}
  \textbf{A.} Final 8-task Avg SR and Forgetting Rate: DTRA~(v2.1) improves Avg SR by $+10.0\%$p ($90.0\%$ vs.\ $80.0\%$) with $0.0\%$ FR.
  \textbf{B.} Inference latency scaling ($K=1 \dots 8$): isolated CPU profiling confirms constant $O(1)$ latency for DTRA~(v2.1) ($3.51 \to 10.98 \to 10.28\,\text{ms}$) versus linear $O(K)$ growth for ABK Blend ($5.08 \to 27.48\,\text{ms}$).
  \textbf{C.} Task magnetic resonance separation margin ($s_{\text{self}} - \max_{j \neq i} s_j$): multimodal keys boost mean margin by $13.9\times$ ($0.2752$ vs.\ $0.0198$).
  \textbf{D.} 3D OOD RGB noise \& brightness perturbation tolerance: DTRA~(v2.1) preserves $89.1\%$ routing accuracy and $42.2\%$ severe OOD SR ($9\times$ higher than ABK's $4.7\%$).}
  \label{fig:mw_benchmark}
\end{figure*}

Table~\ref{tab:mw_results} and Figure~\ref{fig:mw_benchmark} reveal clear gains across every benchmark dimension:
\begin{itemize}[leftmargin=*,itemsep=2pt]
  \item \textbf{Superior Retention on Visually Ambiguous Tasks:} On the difficult later tasks $T_4$--$T_7$ (\texttt{drawer-open/close}, \texttt{handle-press}, \texttt{coffee-button}), ABK Blend drops from $86.7\% \to 53.3\%$ on $T_4$ and $93.3\% \to 73.3\%$ on $T_5$ after $T_5$ and $T_7$ join ($8.57\%$ FR, $-5.71\%$ BWT), and ABK Argmax struggles on $T_4$ ($53.3\%$). DTRA~(v2.1) preserves \textbf{0.0\% FR} while achieving \textbf{93.3\%} on $T_4$, \textbf{100.0\%} on $T_5$, \textbf{100.0\%} on $T_6$, and \textbf{86.7\%} on $T_7$, lifting overall Avg SR to \textbf{90.0\%}.
  \item \textbf{13.9$\times$ Wider Inter-Task Separation:} Table~\ref{tab:efficiency_ood} shows DoG visual prototypes crowd near cosine threshold $0.996$, leaving margins as narrow as $0.0082$ ($T_4$) and $0.0083$ ($T_5$) with mean $0.0198$. DTRA~(v2.1) multimodal keys separate all tasks cleanly (mean margin $0.2752$, up to $0.4575$ on $T_7$, minimum $0.0230$).
  \item \textbf{9$\times$ Higher Severe OOD Tolerance:} Under severe RGB Gaussian noise ($\sigma=0.20$) and brightness shifts ($1.25\times + 0.12$), ABK routing degrades from $100.0\% \to 40.6\%$ and control SR collapses to $4.7\%$ ($0\%$ on $6$ of $8$ tasks). DTRA~(v2.1) sustains \textbf{89.1\%} routing accuracy and \textbf{42.2\%} SR ($75.0\%$ each on \texttt{drawer-close}, \texttt{handle-press}, and \texttt{coffee-button}).
\end{itemize}

\begin{table*}[t]
  \centering
  \caption{Latency scaling ($K=1\dots 8$), magnetic resonance margins, and 3D OOD stress test comparison ($n_{\text{ood}}=8$).}
  \label{tab:efficiency_ood}
  \small
  \resizebox{\textwidth}{!}{
  \begin{tabular}{lcccccccccc}
    \toprule
    & \multicolumn{3}{c}{\textbf{CPU Latency (ms) $\downarrow$}} & \multicolumn{2}{c}{\textbf{Separation Margin $\uparrow$}} & \multicolumn{3}{c}{\textbf{OOD Avg SR (\%) $\uparrow$}} & \multicolumn{2}{c}{\textbf{OOD Route Acc (\%) $\uparrow$}} \\
    \cmidrule(lr){2-4} \cmidrule(lr){5-6} \cmidrule(lr){7-9} \cmidrule(lr){10-11}
    \textbf{Method} & $K=1$ & $K=4$ & $K=8$ & Mean & Min & Clean & Mild & Severe & Clean & Severe \\
    \midrule
    v1.1.3 ABK (Blend) & 5.08 & 17.54 & 27.48 & 0.0198 & 0.0082 & 87.5 & 4.7 & 4.7 & \textbf{100.0} & 40.6 \\
    \textbf{DTRA~(v2.1, Ours)} & \textbf{3.51} & \textbf{10.61} & \textbf{10.28} & \textbf{0.2752} & \textbf{0.0230} & \textbf{90.6} & \textbf{42.2} & \textbf{42.2} & 85.9 & \textbf{89.1} \\
    \bottomrule
  \end{tabular}
  }
\end{table*}

% ==========================================================================
\section{Discussion and Limitations}
% ==========================================================================

\subsection{Key Scan \texorpdfstring{$O(K)$}{O(K)} vs.\ Active-Tube Kernel \texorpdfstring{$O(|A_t|)$}{O(|At|)} and XPU Jitter}
A careful complexity distinction is essential for honest systems analysis. At inference step $t$, scoring cosine similarity between query $\mathbf{x}_t \in \mathbb{R}^{128}$ and the key bank $\mathbf{K} \in \mathbb{R}^{K \times 128}$ entails one matrix--vector product scaling linearly as $O(K D_{\text{key}})$ (\texttt{512\,bytes/task} in \texttt{fp32}). However, all expensive nonlinear operations---closed-form $\operatorname{atan2}$ phase projection, hyperspherical distortion, Soft Tube Overlap, Coarse-to-Fine B-Spline MLP evaluation, and snap-back---run strictly over the sliced active set $A_t$ ($|A_t|=1$ per stage in 3D eval, $1/4$ Coarse \& $2/8$ Fine in Step~1), achieving true \textbf{$O(|A_t|) \approx O(1)$} core computation. Isolated CPU benchmarking verifies this flat scaling ($10.98\,\text{ms}$ at $K=2 \to 10.61\,\text{ms}$ at $K=4 \to 10.28\,\text{ms}$ at $K=8$, a $2.67\times$ speedup over ABK Blend's $27.48\,\text{ms}$). Meanwhile, during interleaved training--evaluation runs on Windows Intel Arc XPU, transient spikes appeared mid-sequence ($27.15$--$37.10\,\text{ms}$ for $K=3 \dots 7$ under DTRA~(v2.1) and $26.07$--$31.67\,\text{ms}$ for $K=3 \dots 5$ under ABK Blend) before converging to $5.43\,\text{ms}$ (DTRA~(v2.1)) and $8.94\,\text{ms}$ (ABK) at $K=8$. This reflects XPU JIT shader compilation and host--device synchronization jitter rather than algorithmic $O(K)$ growth, underscoring why isolated CPU sweeps provide the cleanest test of asymptotic scaling.

\subsection{Stage-1 Primitive Sharing Across Button-Pressing Tasks}
Inspecting per-task routing at $K=8$ on clean held-out episodes ($n_{\text{eval}}=15$) reveals an instructive phenomenon: $6$ of $8$ tasks achieve $100.0\%$ Fine routing accuracy and $T_6$ (\texttt{handle-press-v3}) achieves $80.0\%$ ($100.0\%$ SR), whereas $T_0$ (\texttt{button-press-v3}) records $33.3\%$ Fine top-1 routing accuracy yet sustains \textbf{93.3\% success rate} with \textbf{0.0\% forgetting}. Examination of the held-out similarity matrix shows \texttt{button-press-v3} queries aligning closely with both Tube~$0$ ($0.731$) and Tube~$7$ (\texttt{coffee-button-v3}, $0.708$; margin $0.0230$) because their commands share button-pressing tokens (\emph{``reach''}, \emph{``press''}, \emph{``button''}) and top-down pushing geometry. When Stage~1 Coarse docking dispatches such frames toward shared button-pressing primitives, Stage~2 fine residuals trained on button actuation transfer constructively to complete $T_0$ without interference. Conversely, under OOD noise, preserving $89.1\%$ routing protects clean pushing tasks (\texttt{drawer-close}, \texttt{handle-press}, and \texttt{coffee-button} at $75.0\%$--$87.5\%$ under Mild and $75.0\%$ under Severe noise) while raw Conv4 motor-feature corruption still limits articulated pulling/sliding tasks ($T_1, T_2, T_4$ drop to $0\%$ at $\sigma=0.10$), bounding overall noisy SR to $42.2\%$.

\subsection{Scalability and Future Extensions}
First, extending DTRA beyond two Coarse/Fine stages into a \textbf{Deep Multi-Ring Stack} (e.g., compositional sub-goal sequencing along recurrently myelinated trajectories) will broaden expressivity for long-horizon bimanual manipulation. Second, when task libraries expand to $K \gg 10^3$, replacing the flat key matrix dot product with hierarchical bucketed indexing or Locality-Sensitive Hashing over Coarse clusters will reduce even the initial scan from $O(K)$ to sub-linear $O(\log K)$. Third, substituting hashed unigram/3-gram language projections and Conv4 spatial keys with pretrained vision-language foundation encoders will unlock open-vocabulary physical robot deployment.

% ==========================================================================
\section{Conclusion}
% ==========================================================================

We presented \textbf{DTRA}, modeling continual visuomotor skills as closed first-order elliptical hypertube rings embedded on $\mathbb{S}^{63}$. Uniting multimodal magnetic keys, closed-form $\operatorname{atan2}$ phase projection, bounded elastic compliance, command-gated myelination, two-stage Coarse-to-Fine active-tube docking ($O(|A_t|)\approx O(1)$), and exact post-inference snap-back eliminates catastrophic forgetting ($0.0\%$ FR) while scaling 8-task 3D Meta-World success rate to \textbf{90.0\%} (DTRA~(v2.1)), CPU inference efficiency to \textbf{10.28\,ms} at $K=8$, inter-task resonance margin by \textbf{13.9$\times$}, and severe OOD robustness by \textbf{9$\times$} ($42.2\%$ vs.\ $4.7\%$).

% ==========================================================================
\begin{thebibliography}{15}
\footnotesize
\setlength{\itemsep}{0.9pt}
\setlength{\parskip}{0pt}

\bibitem{son2026abkvla}
S.~Son, ``ABK-VLA: Task-ID-free continual visuomotor policy learning via autonomous B-spline knot compression and visual prototype routing,'' \emph{Zenodo Preprint}, DOI: 10.5281/zenodo.24208545, Nov. 2026.

\bibitem{kim2024openvla}
M.~J. Kim, K.~Pertsch, S.~Karamcheti, et~al., ``OpenVLA: An open-source vision-language-action model,'' \emph{arXiv preprint}, 2024.

\bibitem{zhao2023act}
T.~Z. Zhao, V.~Kumar, S.~Levine, and C.~Finn, ``Learning fine-grained bimanual manipulation with low-cost hardware,'' in \emph{Proc. Robot.: Sci. Syst. (RSS)}, 2023.

\bibitem{chi2023diffusion}
C.~Chi, S.~Feng, Y.~Du, Z.~Xu, E.~Cousineau, B.~Burchfiel, and S.~Song, ``Diffusion policy: Visuomotor policy learning via action diffusion,'' in \emph{Proc. Robot.: Sci. Syst. (RSS)}, 2023.

\bibitem{levine2016visuomotor}
S.~Levine, C.~Finn, T.~Darrell, and P.~Abbeel, ``End-to-end training of deep visuomotor policies,'' \emph{J. Mach. Learn. Res. (JMLR)}, vol.~17, no.~39, pp.~1--40, 2016.

\bibitem{yu2020metaworld}
T.~Yu, D.~Quillen, Z.~He, R.~Julian, A.~Narayan, H.~Shively, A.~Bellathur, K.~Hausman, C.~Finn, and S.~Levine, ``Meta-World: A benchmark and evaluation for multi-task and meta reinforcement learning,'' in \emph{Proc. Conf. Robot Learn. (CoRL)}, 2020.

\bibitem{wolczyk2021continualworld}
M.~Wo{\l}czyk, M.~Zaj\k{a}c, R.~Pascanu, {\L}.~Kuci\'{n}ski, and P.~Mi{\l}o\'{s}, ``Continual World: A robotic benchmark for continual reinforcement learning,'' in \emph{Adv. Neural Inf. Process. Syst. (NeurIPS)}, vol.~34, 2021.

\bibitem{liu2023libero}
B.~Liu, Y.~Zhu, C.~Gao, Y.~Feng, Q.~Liu, Y.~Zhu, and P.~Stone, ``LIBERO: Benchmarking knowledge transfer for lifelong robot learning,'' in \emph{Adv. Neural Inf. Process. Syst. (NeurIPS Datasets and Benchmarks Track)}, 2023.

\bibitem{kirkpatrick2017ewc}
J.~Kirkpatrick et~al., ``Overcoming catastrophic forgetting in neural networks,'' \emph{Proc. Natl. Acad. Sci. (PNAS)}, 2017.

\bibitem{chaudhry2018riemannian}
A.~Chaudhry, P.~K. Dokania, T.~Ajanthan, and P.~H.~S. Torr, ``Riemannian walk for incremental learning: Understanding forgetting and intransigence,'' in \emph{Proc. Eur. Conf. Comput. Vis. (ECCV)}, 2018.

\bibitem{lopezpaz2017gem}
D.~Lopez-Paz and M.~Ranzato, ``Gradient episodic memory for continual learning,'' in \emph{Adv. Neural Inf. Process. Syst. (NeurIPS)}, 2017.

\bibitem{rusu2016progressive}
A.~A. Rusu, N.~C. Rabinowitz, G.~Desjardins, H.~Soyer, J.~Kirkpatrick, K.~Kavukcuoglu, and R.~Hadsell, ``Progressive neural networks,'' \emph{arXiv:1606.04671}, 2016.

\bibitem{mallya2018packnet}
A.~Mallya and S.~Lazebnik, ``PackNet: Adding multiple tasks to a single network by iterative pruning,'' in \emph{Proc. IEEE/CVF Conf. Comput. Vis. Pattern Recognit. (CVPR)}, 2018.

\bibitem{wang2022l2p}
Z.~Wang et~al., ``Learning to prompt for continual learning,'' in \emph{Proc. IEEE/CVF Conf. Comput. Vis. Pattern Recognit. (CVPR)}, 2022.

\bibitem{paraschos2013promp}
A.~Paraschos, C.~Daniel, J.~Peters, and G.~Neumann, ``Probabilistic movement primitives,'' in \emph{Adv. Neural Inf. Process. Syst. (NeurIPS)}, vol.~26, 2013.

\end{thebibliography}

\end{document}
